\documentclass[a4paper,12pt]{article}
\usepackage[T2A]{fontenc}
\usepackage[utf8]{inputenc}
\usepackage[russian]{babel}
\usepackage{amsmath, amssymb}
\usepackage{geometry}
\geometry{top=2cm, bottom=2cm, left=2.5cm, right=1.5cm}
\usepackage{enumitem}
\setlist[enumerate]{leftmargin=*, itemsep=3pt, parsep=0pt}

\title{Ответы на вопросы по интерполяции}
\author{}
\date{}

\begin{document}

\maketitle

\begin{enumerate}

\item \textbf{Когда возникает необходимость в использовании интерполяционных методов?}

	Интерполяционные методы применяются в тех случаях, когда функция \( y = f(x) \) задана не аналитической формулой, а только в дискретных точках (узлах интерполяции), но требуется найти её значение в промежуточных точках, не совпадающих с узлами. Например, при обработке экспериментальных данных, в численном анализе, в задачах компьютерной графики, при аппроксимации сложных функций более простыми и т.д.

\item \textbf{Чем отличается аппроксимация от интерполяции?}

	\textit{Интерполяция} требует точного прохождения построенной функции через все заданные узловые точки. \textit{Аппроксимация} (например, методом наименьших квадратов) ищет функцию, которая приближает исходные данные в некотором интегральном смысле, и не обязана проходить через узлы. Аппроксимация используется, когда данные содержат погрешности (шумы).

\item \textbf{В чём сущность задачи интерполирования?}

	Сущность задачи интерполирования: для функции, заданной таблично \( (x_i, y_i), i=0,\dots,n \), построить более простую функцию \( P(x) \) (например, многочлен), такую что \( P(x_i) = y_i \) для всех \( i \), и затем использовать \( P(x) \) для вычисления значений в произвольных точках \( x \in [x_0, x_n] \).

\item \textbf{Поясните смысл терминов: интерполяция, экстраполяция.}

	\textit{Интерполяция} — нахождение значений функции внутри интервала, содержащего узлы \( [x_0, x_n] \). \textit{Экстраполяция} — нахождение значений за пределами этого интервала (левее \( x_0 \) или правее \( x_n \)). Экстраполяция значительно менее точна и надёжна.

\item \textbf{Как найти приближенное значение функции при линейной интерполяции?}

	Линейная интерполяция использует две соседние точки \( (x_0, y_0) \) и \( (x_1, y_1) \):
	\[
	P_1(x) = y_0 + \frac{y_1 - y_0}{x_1 - x_0} (x - x_0).
	\]
	Для \( x \in [x_0, x_1] \) подставляем значение и вычисляем.

\item \textbf{Как найти приближенное значение функции при квадратичной интерполяции?}

	Квадратичная интерполяция использует три точки \( (x_0, y_0), (x_1, y_1), (x_2, y_2) \). Строится многочлен второй степени:
	\[
	P_2(x) = a_0 + a_1 (x - x_0) + a_2 (x - x_0)(x - x_1),
	\]
	где коэффициенты находятся из условий \( P_2(x_i) = y_i \). Затем подставляем нужное \( x \).

\item \textbf{Как строится интерполяционный многочлен Лагранжа?}

	Многочлен Лагранжа степени \( n \) имеет вид:
	\[
	L_n(x) = \sum_{i=0}^{n} y_i \cdot \ell_i(x), \quad
	\ell_i(x) = \prod_{\substack{j=0 \\ j \neq i}}^{n} \frac{x - x_j}{x_i - x_j}.
	\]
	Здесь \( \ell_i(x_i) = 1 \) и \( \ell_i(x_j) = 0 \) при \( j \neq i \), поэтому \( L_n(x_i) = y_i \).

\item \textbf{Дайте определение понятий разделенной разности нулевого и первого порядков.}

	\begin{itemize}
		\item Разделённая разность нулевого порядка: \( f[x_i] = y_i \).
		\item Разделённая разность первого порядка:
		\[
		f[x_i, x_{i+1}] = \frac{f[x_{i+1}] - f[x_i]}{x_{i+1} - x_i}.
		\]
		Она характеризует среднюю скорость изменения функции на отрезке.
	\end{itemize}

\item \textbf{Объясните принцип построения интерполяционного полинома Ньютона.}

	Полином Ньютона с разделёнными разностями строится как:
	\[
	P_n(x) = f[x_0] + f[x_0,x_1](x-x_0) + f[x_0,x_1,x_2](x-x_0)(x-x_1) + \dots
	\]
	Коэффициенты — разделённые разности. Преимущество: при добавлении нового узла можно пересчитать только последние слагаемые (рекуррентность).

\item \textbf{Покажите графическую интерпретацию интерполяции.}

	На графике: заданы точки \( (x_i, y_i) \). Интерполяция — построение кривой (многочлена), проходящей через все эти точки. Визуально — гладкая линия, соединяющая отметки. Значение в промежуточной точке снимается с этой линии.

\item \textbf{В каких случаях используются конечные разности, в каких – разделенные?}

	\textit{Конечные разности} применяются при равномерном шаге узлов (\( h = x_{i+1} - x_i = const \)). \textit{Разделённые разности} — при произвольном (неравномерном) расположении узлов.

\item \textbf{В каких случаях используют формулу Ньютона для интерполирования вперед и для интерполирования назад?}

	\begin{itemize}
		\item \textbf{Вперёд}: когда аргумент \( x \) находится в начале таблицы (\( x \) близок к \( x_0 \)). Используются конечные разности вверх по столбцу.
		\item \textbf{Назад}: когда \( x \) находится в конце таблицы (\( x \) близок к \( x_n \)). Используются конечные разности вниз по столбцу.
	\end{itemize}

\item \textbf{В каких случаях используют формулу Гаусса для интерполирования вперед и для интерполирования назад?}

	Формулы Гаусса применяются при равномерном шаге для интерполяции вблизи середины таблицы. \textit{Первая формула Гаусса} (вперёд) использует узлы \( x_0, x_1, x_{-1}, x_2, \dots \). \textit{Вторая формула Гаусса} (назад) — \( x_0, x_{-1}, x_1, x_{-2}, \dots \). Выбор зависит от того, с какой стороны от \( x_0 \) находится искомая точка.

\item \textbf{В каких случаях используют формулу Стирлинга?}

	Формула Стирлинга применяется при равномерном шаге для интерполяции вблизи центра таблицы, когда используются симметрично расположенные узлы. Она является средним арифметическим двух формул Гаусса и даёт более высокую точность для точек около \( x_0 \).

\item \textbf{В каких случаях используют формулу Бесселя?}

	Формула Бесселя также применяется для центральной интерполяции при равномерном шаге, но она ориентирована на середину между узлами. Используется, когда \( x \) находится вблизи \( x_0 + \frac{h}{2} \). Выражается через полуразности конечных разностей.

\item \textbf{В чем разница между глобальной и локальной разновидностями интерполяции?}

	\begin{itemize}
		\item \textit{Глобальная интерполяция}: один многочлен высокой степени на всём интервале. Плюс: гладкость. Минус: неустойчивость, сильные осцилляции (феномен Рунге).
		\item \textit{Локальная интерполяция}: на каждом участке строится свой многочлен невысокой степени (например, кусочная линейная или кубическая интерполяция). Плюс: устойчивость. Минус: потеря гладкости в стыках (кроме сплайнов).
	\end{itemize}

\item \textbf{В чем заключается интерполяция кубическими сплайнами?}

	Кубический сплайн — это кусочно-полиномиальная функция третьей степени, непрерывная вместе с первой и второй производными в узлах сетки. На каждом промежутке \( [x_i, x_{i+1}] \) строится свой кубический многочлен, и они сшиваются гладко. Это даёт хорошее приближение без осцилляций, особенно для плавных функций.

\end{enumerate}

\end{document}